% Chapter 4, Topic _Linear Algebra_ Jim Hefferon
%  http://joshua.smcvt.edu/linearalgebra
%  2001-Jun-12
\topic{Metode Pangkat}
\index{metode pangkat|(}
\index{pangkat!metode|(}
Dalam penerapan, matriks dapat berukuran besar. Menghitung nilai eigen dan
vektor eigen dengan mencari lalu menyelesaikan polinomial karakteristik tidak
praktis, terlalu lambat, dan terlalu rentan kesalahan. Beberapa teknik
menghindari polinomial karakteristik. Di sini kita membahas suatu metode yang
cocok untuk matriks besar yang \definend{jarang}\index{matriks!jarang}%
\index{matriks jarang}, yakni matriks yang sebagian besar entrinya nol.

Andaikan matriks $T$ berukuran $\nbyn{n}$ memiliki $n$ nilai eigen berbeda
$\lambda_1$, $\lambda_2$, \ldots, $\lambda_n$. Maka $\C^n$ memiliki basis yang
terdiri atas vektor-vektor eigen terkait
$\sequence{\vec{\zeta}_1,\dots,\vec{\zeta}_n}$. Untuk sebarang
$\vec{v}\in\C^n$, tulis
$\vec{v}=c_1\vec{\zeta}_1+\dots+c_n\vec{\zeta}_n$ dan iterasikan $T$ pada
$\vec{v}$ untuk memperoleh
\begin{align*}
  T\vec{v} 
      &=c_1\lambda_1\vec{\zeta}_1+c_2\lambda_2\vec{\zeta}_2+
                              \dots+c_n\lambda_n\vec{\zeta}_n  \\
  T^2\vec{v} 
      &=c_1\lambda_1^2\vec{\zeta}_1+c_2\lambda_2^2\vec{\zeta}_2+
                              \dots+c_n\lambda_n^2\vec{\zeta}_n  \\
  T^3\vec{v} 
      &=c_1\lambda_1^3\vec{\zeta}_1+c_2\lambda_2^3\vec{\zeta}_2+
                              \dots+c_n\lambda_n^3\vec{\zeta}_n  \\
      &\vdotswithin{=}                                            \\
  T^k\vec{v} 
      &=c_1\lambda_1^k\vec{\zeta}_1+c_2\lambda_2^k\vec{\zeta}_2+
                              \dots+c_n\lambda_n^k\vec{\zeta}_n  
\end{align*}
Dengan mengandaikan bahwa $|\lambda_1|$ benar-benar lebih besar daripada
semua $|\lambda_i|$ lainnya dan membagi dengan $\lambda_1^k$, diperoleh
\begin{equation*}
  \frac{T^k\vec{v}}{\lambda_1^k} 
      =c_1\vec{\zeta}_1+c_2\frac{\lambda_2^k}{\lambda_1^k}\vec{\zeta}_2+
                    \dots+c_n\frac{\lambda_n^k}{\lambda_1^k}\vec{\zeta}_n  
\end{equation*} 
Ketika $k$ membesar, pecahan-pecahan tersebut menuju nol. Jadi, suku milik
$\lambda_1$ mendominasi ekspresi itu dan ekspresinya menuju
$c_1\vec{\zeta}_1$.

Jadi, jika $c_1\neq 0$, ketika $k$ membesar, arah vektor $T^k\vec{v}$
mendekati arah vektor eigen yang terkait dengan nilai eigen dominan.
Akibatnya, rasio panjang vektor
$\absval{T^{k}\vec{v}}/\absval{T^{k-1}\vec{v}}$ mendekati nilai mutlak nilai
eigen dominan tersebut.

Sebagai contoh, nilai eigen matriks 
\begin{equation*}
  T=\begin{mat}[r]
    3  &0  \\
    8  &-1
  \end{mat}
\end{equation*}
adalah $3$ dan $-1$. Jika komponen $\vec{v}$ adalah $1$ dan $1$, iterasi
memberikan hasil berikut.
\begin{center}
  \begin{tabular}{c|ccccc}
     $\vec{v}$  &$T\vec{v}$  &$T^2\vec{v}$ 
        &$\cdots$ &$T^9\vec{v}$ &$T^{10}\vec{v}$        \\ \hline 
     $\matrixvenlarge{\colvec[r]{1 \\ 1}}$  
        &$\matrixvenlarge{\colvec[r]{3 \\ 7}}$ 
        &$\matrixvenlarge{\colvec[r]{9 \\ 17}}$ 
        &$\cdots$  
        &$\matrixvenlarge{\colvec[r]{19\,683 \\ 39\,367}}$   
        &$\matrixvenlarge{\colvec[r]{59\,049 \\ 118\,097}}$
  \end{tabular}
\end{center}
Rasio panjang dua vektor terakhir adalah $2.999\,9$.

Ada dua persoalan implementasi. Pertama, alih-alih menentukan pangkat-pangkat
$T$ lalu menerapkannya pada $\vec{v}$, kita hitung $\vec{v}_1=T\vec{v}$,
kemudian $\vec{v}_2=T\vec{v}_1$, dan seterusnya (jadi, $T^2\!$, $T^3\!$,
\ldots{} tidak dihitung secara terpisah). Hasil kali matriks--vektor ini dapat
dihitung dengan cepat meskipun $T$ besar, asalkan $T$ jarang.

Kedua, agar tidak menghasilkan bilangan yang begitu besar hingga melampaui
kemampuan komputer, kita dapat menormalkan $\vec{v}_i$ pada setiap langkah.
Misalnya, bagi setiap $\vec{v}_i$ dengan panjangnya (pilihan lain adalah
membaginya dengan komponen terbesar atau cukup dengan komponen pertamanya).
Dengan demikian, metode ini diimplementasikan dengan menghasilkan
\begin{align*}
  \vec{w}_0  &=\vec{v}_0/\absval{\vec{v}_0} \\
  \vec{v}_1  &=T\vec{w}_0                 \\
  \vec{w}_1  &=\vec{v}_1/\absval{\vec{v}_1} \\
  \vec{v}_2  &=T\vec{w}_1                 \\
             &\vdotswithin{=}    \\
  \vec{w}_{k-1}  &=\vec{v}_{k-1}/\absval{\vec{v}_{k-1}} \\
  \vec{v}_k  &=T\vec{w}_{k-1}                 
\end{align*}
hingga pendekatannya dianggap memadai. Kemudian $\vec{v}_k$ mendekati suatu
vektor eigen, dan nilai eigen dominannya dapat didekati dengan hasil bagi
Rayleigh % $\absval{\vec{v}_{k}}/\absval{\vec{w}_{k-1}}$.
$((T\vec{v}_k)\dotprod \vec{v}_k)/(\vec{v}_k\dotprod\vec{v}_k)
  \approx (\lambda_1\vec{v}_k\cdot\vec{v}_k)/(\vec{v}_k\dotprod\vec{v}_k)
  =\lambda_1$.


Salah satu kriteria kecukupan adalah mengiterasikan hingga pendekatan nilai
eigen menjadi stabil. Misalnya, proses iterasi dapat dihentikan bukan setelah
sejumlah langkah tetap, melainkan ketika dua pendekatan hasil bagi Rayleigh
berturut-turut berbeda kurang dari satu persen atau sama hingga digit signifikan
kedua.

Laju konvergensi ditentukan oleh laju pangkat-pangkat
$|\lambda_2/\lambda_1|$ menuju nol, dengan $\lambda_2$ sebagai nilai eigen
yang nilai mutlaknya terbesar kedua. Jika rasio itu jauh lebih kecil daripada
satu, konvergensinya cepat; jika hanya sedikit lebih kecil daripada satu,
konvergensinya dapat sangat lambat. Karena itu, metode pangkat bukan cara yang
paling umum digunakan untuk menentukan nilai eigen (meskipun merupakan cara
yang paling sederhana, sehingga dibahas di sini). Berbagai metode lain biasanya
mula-mula mengganti matriks $T$ dengan matriks serupa yang memiliki nilai eigen
sama, tetapi berada dalam bentuk tereduksi seperti
\definend{bentuk tridiagonal}\index{bentuk tridiagonal}%
\index{matriks!tridiagonal}, yang hanya memiliki entri tak nol pada diagonal
atau tepat di atas atau di bawahnya. Teknik-teknik kasus khusus kemudian dapat
menentukan nilai eigennya. Setelah nilai eigen diketahui, vektor eigen $T$
dapat dihitung dengan mudah. Metode-metode lain itu berada di luar cakupan kita.
Rujukan yang baik adalah \cite{Goult}.



\begin{exercises}
  \item 
    Gunakan sepuluh iterasi untuk memperkirakan nilai eigen dominan matriks-
    matriks berikut, dimulai dari vektor berkomponen $1$ dan $2$.
    Bandingkan jawabannya dengan hasil penyelesaian persamaan karakteristik.
    \begin{exparts*}
      \partsitem $\begin{mat}[r]
                    1  &5  \\
                    0  &4
                  \end{mat}$
      \partsitem $\begin{mat}[r]
                    3   &2  \\
                    -1  &0
                  \end{mat}$
    \end{exparts*}
    \begin{answer}
     \begin{exparts}
       \partsitem Dengan pengamatan langsung, nilai eigen dominannya adalah $4$.
         \Sage{} memberikan hasil berikut.
\begin{lstlisting}
sage: def eigen(M,v,num_loops=10):
....:     for p in range(num_loops):
....:         v_normalized = (1/v.norm())*v
....:         v = M*v_normalized
....:     return v 
....: 
sage: M = matrix(RDF, [[1,5], [0,4]])
sage: v = vector(RDF, [1, 2])
sage: v = eigen(M,v)
sage: (M*v).dot_product(v)/v.dot_product(v)
4.00000147259
\end{lstlisting}
       \partsitem Perhitungan sederhana menunjukkan bahwa nilai eigen dominannya
          adalah~$2$.
          \Sage{} memberikan hasil berikut.
\begin{lstlisting}
sage: M = matrix(RDF, [[3,2], [-1,0]])
sage: v = vector(RDF, [1, 2])
sage: v = eigen(M,v)
sage: (M*v).dot_product(v)/v.dot_product(v)
2.00097741083
\end{lstlisting}
     \end{exparts}
    \end{answer}
  \item 
     Ulangi latihan sebelumnya dengan mengiterasikan hingga jarak antara dua
     vektor ternormalisasi berturut-turut kurang dari $0.01$.
     Pada setiap langkah, normalkan dengan membagi setiap vektor dengan
     panjangnya. Berapa iterasi yang diperlukan? Apakah jawabannya berbeda
     secara berarti?
     \begin{answer}
       \begin{exparts}
         \partsitem Berikut hasil dari \Sage{}.
 \begin{lstlisting}
sage: def eigen_by_iter(M, v, toler=0.01):
....:     dex = 0
....:     diff = 10
....:     while abs(diff)>toler:
....:         dex = dex+1
....:         v_next = M*v
....:         v_normalized = (1/v.norm())*v
....:         v_next_normalized = (1/v_next.norm())*v_next
....:         diff = (v_next_normalized-v_normalized).norm()
....:         v_prior = v_normalized
....:         v = v_next_normalized
....:     return v, v_prior, dex
....: 
sage: M = matrix(RDF, [[1,5], [0,4]])
sage: v = vector(RDF, [1, 2])
sage: v,v_prior,dex = eigen_by_iter(M,v)
sage: (M*v).norm()/v.norm()
4.00604111686
sage: dex
4           
 \end{lstlisting}
      \partsitem \Sage{} memerlukan beberapa iterasi tambahan untuk kasus ini.
        Perhitungan berikut menggunakan prosedur yang didefinisikan pada bagian
        sebelumnya.
\begin{lstlisting}
sage: M = matrix(RDF, [[3,2], [-1,0]])
sage: v = vector(RDF, [1, 2])
sage: v,v_prior,dex = eigen_by_iter(M,v)
sage: (M*v).norm()/v.norm()
2.01585174302
sage: dex
6        
\end{lstlisting}
       \end{exparts}
     \end{answer}
  \item 
    Gunakan sepuluh iterasi untuk memperkirakan nilai eigen dominan matriks-
    matriks berikut, dimulai dari vektor berkomponen $1$, $2$, dan $3$.
    Bandingkan jawabannya dengan hasil penyelesaian persamaan karakteristik.
    \begin{exparts*}
      \partsitem $\begin{mat}[r]
                    4   &0  &1 \\
                    -2  &1  &0  \\
                    -2  &0  &1
                  \end{mat}$
      \partsitem $\begin{mat}[r]
                   -1  &2  &2  \\
                    2  &2  &2  \\
                   -3  &-6 &-6
                  \end{mat}$
    \end{exparts*}
    \begin{answer}
     \begin{exparts}
       \partsitem Nilai eigen dominannya adalah $3$.
         \Sage{} memberikan hasil berikut.
\begin{lstlisting}
sage: M = matrix(RDF, [[4,0,1], [-2,1,0], [-2,0,1]])
sage: v = vector(RDF, [1, 2, 3])
sage: v = eigen(M,v)
sage: (M*v).dot_product(v)/v.dot_product(v)
3.02362112326
\end{lstlisting}
       \partsitem Nilai eigen dominannya adalah $-3$.
\begin{lstlisting}
sage: M = matrix(RDF, [[-1,2,2], [2,2,2], [-3,-6,-6]])
sage: v = vector(RDF, [1, 2, 3])
sage: v = eigen(M,v)
sage: (M*v).dot_product(v)/v.dot_product(v)
-3.00941127145
\end{lstlisting}
     \end{exparts}
    \end{answer}
  \item 
     Ulangi latihan sebelumnya dengan mengiterasikan hingga jarak antara dua
     vektor ternormalisasi berturut-turut kurang dari $0.01$.
     Pada setiap langkah, normalkan dengan membagi setiap vektor dengan
     panjangnya. Berapa iterasi yang diperlukan? Apakah jawabannya berbeda
     secara berarti?
     \begin{answer}
       \begin{exparts}
         \partsitem
           \Sage{} memberikan hasil berikut, dengan
           \lstinline!eigen_by_iter! sebagaimana didefinisikan di atas.
\begin{lstlisting}
sage: M = matrix(RDF, [[4,0,1], [-2,1,0], [-2,0,1]])
sage: v = vector(RDF, [1, 2, 3])
sage: v,v_prior,dex = eigen_by_iter(M,v) 
sage: (M*v).dot_product(v)/v.dot_product(v)
3.05460392934
sage: dex
8            
\end{lstlisting}
         \partsitem Dengan pengaturan berikut,
\begin{lstlisting}
sage: M = matrix(RDF, [[-1,2,2], [2,2,2], [-3,-6,-6]])
sage: v = vector(RDF, [1, 2, 3])       
\end{lstlisting}
            \Sage{} tidak berhenti (gunakan \lstinline!<Ctrl>-c! untuk
            menghentikan perhitungan). Dengan menambahkan pemeriksaan galat
            pada rutin tersebut,
\begin{lstlisting}
def eigen_by_iter(M, v, toler=0.01):
    dex = 0
    diff = 10
    while abs(diff)>toler:
        dex = dex+1
        if dex>1000:
            print("waduh! tampaknya terjadi perulangan: \nv=",v,"\nv_next=",v_next)
        v_next = M*v
        if (v.norm()==0):
            print("waduh! v bernilai nol")
            return None
        if (v_next.norm()==0):
            print("waduh! v_next bernilai nol")
            return None
        v_normalized = (1/v.norm())*v
        v_next_normalized = (1/v_next.norm())*v_next
        diff = (v_next_normalized-v_normalized).norm()
        v_prior = v_normalized
        v = v_next_normalized
    return v, v_prior, dex              
\end{lstlisting}
           diperoleh hasil berikut.
\begin{lstlisting}
waduh! tampaknya terjadi perulangan: 
  v= (0.707106781187, -1.48029736617e-16, -0.707106781187) 
  v_next= (2.12132034356, -4.4408920985e-16, -2.12132034356) 
waduh! tampaknya terjadi perulangan: 
  v= (-0.707106781187, 1.48029736617e-16, 0.707106781187) 
  v_next= (-2.12132034356, 4.4408920985e-16, 2.12132034356)
waduh! tampaknya terjadi perulangan: 
  v= (0.707106781187, -1.48029736617e-16, -0.707106781187) 
  v_next= (2.12132034356, -4.4408920985e-16, -2.12132034356)           
\end{lstlisting}
          Jadi, iterasinya berputar di antara dua arah yang berlawanan.
       \end{exparts}
     \end{answer}
  \item 
      Apa yang terjadi jika $c_1=0$? Artinya, apa yang terjadi jika vektor awal
      tidak memiliki komponen dalam arah vektor eigen yang relevan?
     \begin{answer}
       Secara teori, metode ini akan menghasilkan $\lambda_2$. Namun, dalam
       praktik, galat pembulatan dalam perhitungan memasukkan komponen dalam
       arah $\vec{v}_1$. Karena itu, metode tersebut tetap akan menghasilkan
       $\lambda_1$, meskipun mungkin memerlukan waktu lebih lama dibandingkan
       pilihan vektor awal yang lebih menguntungkan.
     \end{answer}
  \item 
    Bagaimana metode pangkat dapat diadaptasi untuk menentukan nilai eigen
    terkecil dalam nilai mutlak?
    \begin{answer}
      Alih-alih menggunakan $\vec{v}_k=T\vec{v}_{k-1}$, gunakan
      $\vec{v}_k=T^{-1}\vec{v}_{k-1}$, atau secara ekuivalen selesaikan
      $T\vec{v}_k=\vec{v}_{k-1}$ pada setiap langkah, dengan syarat $T$
      nonsingular.
    \end{answer}
\end{exercises}

\announcecomputercode
Berikut kode untuk sistem aljabar komputer Octave yang melakukan perhitungan di
atas. (Kode ini disunting ringan untuk menghapus baris kosong dan sebagainya.)
\begin{lstlisting}
>T=[3, 0;
    8, -1]
T=
   3   0
   8  -1
>v0=[1; 1]
v0=
   1
   1
>v1=T*v0
v1=
   3
   7
>v2=T*v1
v2=
   9
  17
>T9=T**9
T9=
  19683  0
  39368 -1
>T10=T**10
T10=
  59049  0
 118096  1
>v9=T9*v0
v9=
  19683
  39367
>v10=T10*v0
v10=
  59049
 118097
>norm(v10)/norm(v9)
ans=2.9999
\end{lstlisting}
\textit{Catatan.}
Kode ini belum menggunakan seluruh kemampuan Octave; Octave memiliki fungsi
bawaan yang secara otomatis menerapkan metode canggih untuk menentukan nilai
eigen dan vektor eigen.

\index{pangkat!metode|)}
\index{metode pangkat|)}
\endinput

